\cotefile{62}
\scanpage{36}
\begingroup\scriptsize
\diagram{\begin{tikzcd}[ampersand replacement=\&,column sep=.35em,row sep=2.6em,cramped]
t\ P^1\arrow[d,"\mathbb Z/3"',"f"]\arrow[d,phantom,"\scriptstyle\deleted{\unsure{dt''/t}}"',near end]
\&\substack{\underbrace{0,1,\infty}\\\deleted{\unsure{0,1,\infty;\,2,-1,1/2}}}\arrow[d,no head]
\&\substack{\underbrace{2,-1,1/2}\\\deleted{\unsure{2,-1,1/2}}}\arrow[d,no head]
\&\substack{x'\deleted{\unsure{x''}}=\\\tfrac12(1+\sqrt{-3})}\arrow[d,no head,shift left=.7]\arrow[d,no head,"\text{ramifié}"]\arrow[d,no head,shift right=.7]
\&{}\&\substack{x''=\\\tfrac12(1-\sqrt{-3})}\arrow[d,no head,shift left=.7]\arrow[d,no head,"\text{ramifié}"]\arrow[d,no head,shift right=.7]\\
s\ P^1\arrow[d,"\mathbb Z/2"',"g"]
\&\infty\arrow[d,no head,shift left=.5,"\text{ramifié}"]\arrow[d,no head,shift right=.5]
\&0\arrow[d,no head,shift left=.5,"\text{ramifié}"]\arrow[d,no head,shift right=.5]
\&1\deleted{\unsure{x}}\arrow[dr,no head,postaction={decorate,decoration={zigzag,segment length=4pt,amplitude=1.55pt}},"\deleted{\diagdown}"']
\&{\underbrace{\hspace{2em}}}\arrow[d,no head,"\text{non ramifié}"]
\&-1\arrow[dl,no head,postaction={decorate,decoration={zigzag,segment length=5pt,amplitude=1.1pt}},"\deleted{\diagup}"]\\
r\ P^1\&\infty\&0\&{}\&1\&{}
\arrow[from=1-2,to=2-3,phantom,pos=.58,"\substack{\text{non ramifié}\\\swarrow\quad\searrow}" description]
\end{tikzcd}}
\endgroup

Pour déterminer explicitement $f$ et $g$, on procède ainsi. [On \deleted{\illegible{1}} peut supposer que \deleted{$S=\operatorname{Spec}(\mathbb Z\deleted{\unsure{a}}[2^{-1},3^{-1}])$} \inserted{$S=\operatorname{Spec}(A)$, $A=\mathbb Z[2^{-1},3^{-1}]$} puisque \unsure{de fait|de même|de tte|toute} la situation provient par image inverse d'une situation sur $\operatorname{Spec}(A)$]. La restriction de $g$ \unsure{:|à} $P^1_S-(\infty)-(0)$ est alors un \deleted{\unsure{$S$-sch}} morphisme de $G_m$ dans lui-=, donc est de la forme $t\mathbin{\unsure{-|\to|\mapsto}}\lambda t^n$, où $n\in\mathbb Z$, et $\lambda\in\mathrm H^0(S,\underline{\mathcal O}_S^{+})$. Comme $g(\infty)=\infty$, on a $n>0$, et comme $g$ est \inserted{un rev.} d'ordre $2$, on a $n=2$. D'autre part \deleted{$\lambda$ \unsure{est une}} \inserted{$g(1)=1$, d'où $\lambda=1$}

\deleted{unité de $A$, donc de la forme $2^\alpha 3^\beta$.}

\begingroup\csname tagsleft@true\endcsname
\[
r=g(s)=s^2 \tag*{(9)}
\]
\endgroup

\uline{Cette formule est valable sur tout $P^1$} [lui donner un sens, et justifier cette assertion --- ]
% archive sheet number: 35
